\documentclass[lettersize,journal]{IEEEtran}
\usepackage{amsmath,amsfonts}
\usepackage{algorithmic}
\usepackage{algorithm}
\usepackage{array}
\usepackage[caption=false,font=normalsize,labelfont=sf,textfont=sf]{subfig}
\usepackage{textcomp}
\usepackage{stfloats}
\usepackage{url}
\usepackage{verbatim}
\usepackage{graphicx}
\usepackage{cite}
\usepackage{import}
\hyphenation{op-tical net-works semi-conduc-tor IEEE-Xplore}
% updated with editorial comments 8/9/2021

\usepackage{color}
\usepackage{xcolor}
\usepackage{multirow}
\usepackage{colortbl}
\usepackage{picinpar}
\usepackage{stfloats}
\usepackage{url}
% \usepackage{algorithm,float,algpseudocode} % algpseudocode

\usepackage{mathtools}
\usepackage{listings}                 % Source code printer for LaTeX
% \usepackage{booktabs}
\usepackage{balance}
% \usepackage[ruled, vlined, linesnumbered]{algorithm2e}
% \usepackage{algpseudocode}
\usepackage{ragged2e} 
\usepackage{booktabs}
\usepackage{multirow}
\usepackage{threeparttable}
% \usepackage{subfigure}

% % \usepackage[hyphens]{url}             % Allow hyphens in URL's
% % \usepackage[unicode=false,psdextra,citecolor=green]{hyperref}                 % References package
\usepackage[pagebackref=true,breaklinks=true,colorlinks,bookmarks=false]{hyperref}
% % \modulolinenumbers[5]

\usepackage[noabbrev,nameinlink,capitalise]{cleveref} % Clever references. Options: "fig. !1!" --> "!Figure 1!"
\Crefformat{figure}{#2Fig.~#1#3}
\Crefmultiformat{figure}{Figs.~#2#1#3}{ and~#2#1#3}{, #2#1#3}{ and~#2#1#3}

\DeclareMathOperator*{\argminA}{arg\,min} % Jan Hlavacek
\DeclareMathOperator*{\argminB}{argmin}   % Jan Hlavacek
\newcommand{\mbf}[1]{\ensuremath{\mathbf{#1}}}
\newcommand{\mdim}[2]{\ensuremath{#1\times#2}}
\newcommand{\minv}[1]{\ensuremath{\mathbf{#1}^{-1}}}
\newcommand{\mtrans}[1]{\ensuremath{\mathbf{#1}^{\top}}}
\newcommand{\mbft}[2]{\ensuremath{\mathbf{#1}_\text{#2}}}
\newcommand{\mbfttrans}[2]{\ensuremath{\mathbf{#1}^\top_\text{#2}}}
\newcommand{\mb}[1]{\mathbf{#1}}
\newcommand{\mf}[1]{\mathcal{F}_{#1}}

% 设置 hyperref 的颜色
\hypersetup{
    colorlinks=true, % 启用颜色链接
    linkcolor=red, % 内部链接颜色
    citecolor=green, % 引用颜色
    urlcolor=red, % URL 颜色
    linkbordercolor=red, % 内部链接边框颜色
    citebordercolor=green, % 引用边框颜色
    urlbordercolor=red % URL 边框颜色
}

% \newcommand{\frame}[1]{${\mathcal{#1}}$}
\newcommand{\coor}[1]{#1 coordinate frame}
% \DeclareRobustCommand{\mbb}[1]{\mathbb{#1}}
\DeclareRobustCommand{\mso}{\mathfrak{so}}
\DeclareRobustCommand{\mse}{\mathfrak{se}}

%%%%%很重要，完稿之后注释下面两行。
% \usepackage{CJKutf8}
\setlength{\marginparwidth}{2cm}
\usepackage{todonotes}

\begin{document}

\title{
A Multi-Sensor Measurement System for Urban Altitude Estimation Using Physics-Informed Neural Fields
}

\author{
        Zichun~Chen$^{1}$,
        Siyu~Zhu$^{2}$,
        Shuo~Sun$^{1}$,
        Yanxing~Wu$^{1}$,
        Yongxing~Wang$^{1}$,
        and~Xiao~Hu$^{1}$
\thanks{Corresponding author: Xiao Hu.}
\thanks{$^{1}$Lower AirSpace Economy Research Institute, International Digital Economy Academy, Shenzhen 510085, China (e-mail: chenzichun@idea.edu.cn, sunshuo@idea.edu.cn, wuyanxing@idea.edu.cn, wangyongxing@idea.edu.cn, huxiao1@idea.edu.cn).}
\thanks{$^{2}$the School of Microelectronics, Xi'an Jiaotong University, Xi'an 710049, Shaanxi Province, China (e-mail: zhusiyu@stu.xjtu.edu.cn).
}
        % <-this % stops a space
% \thanks{This paper was produced by the IEEE Publication Technology Group. They are in Piscataway, NJ.}% <-this % stops a space
% \thanks{Manuscript received April 19, 2021; revised August 16, 2021.}
}

% The paper headers
\markboth{Submission to IEEE Transactions on Instrumentation and Measurement}%
{Shell \MakeLowercase{\textit{et al.}}: A Sample Article Using IEEEtran.cls for IEEE Journals}

% \IEEEpubid{0000--0000/00\$00.00~\copyright~2021 IEEE}
% Remember, if you use this you must call \IEEEpubidadjcol in the second
% column for its text to clear the IEEEpubid mark.

\maketitle

\begin{abstract}
Accurate vertical positioning in low airspace is a critical challenge for the safety of urban air mobility and the integration of heterogeneous drone operations. While barometric altimetry provides a baseline via the hydrostatic equation, microclimate variations, urban canyon effects, and low-cost sensor biases typically introduce severe altitude discrepancies. To address these measurement challenges, this paper proposes a novel multi-sensor fusion framework comprising a custom low-cost edge-sensing device and an advanced physics-informed neural field algorithm. The proposed approach integrates physics-based atmospheric constraints with multi-resolution hash encoding, terrain-aware feature engineering, and a curriculum learning strategy to reconstruct highly accurate urban altitude fields from sparse IoT measurements. Validated using a real-world dataset, the system demonstrates exceptional measurement reliability. Experimental results demonstrate that coupling purpose-built IoT instrumentation with spatial-neural architectures can reliably achieve meter-level accuracy, providing a robust, real-time unified height reference for emerging low-altitude economies.
\end{abstract}

\begin{IEEEkeywords}
Neural Fields, Altitude Estimation, Urban Sensing, Curriculum Learning, Hash Encoding, Barometric Pressure.
\end{IEEEkeywords}



% --- Teaser Image Start ---
\begin{figure*}[!ht]
\centering
\includegraphics[width=0.9\textwidth]{image/teaser1.pdf} % 请替换您的图片文件名
\caption{The proposed edge-sensing and altitude harmonization framework. (a) Heterogeneous vertical navigation systems (Barometric vs. GNSS MSL) create severe altitude discrepancies and collision risks in VLL urban airspace. (b) Distributed GeoBox edge terminals are deployed to collect sparse, synchronized 4D atmospheric and geometric data. (c) The proposed PINF algorithm reconstructs these sparse measurements into a continuous, high-precision unified geometric height transformation field for safe U-space operations.}
    \label{fig:concept}
\end{figure*}
% --- Teaser Image End ---
\section{Introduction}
\label{sec:intro}

\IEEEPARstart{T}{he} rapid proliferation of Unmanned Aircraft Systems (UAS) and the emerging Urban Air Mobility (UAM) sector have fundamentally transformed the operational paradigm of Very Low-Level (VLL) airspace. A critical, yet unresolved, challenge in ensuring the safety of high-density air traffic is the harmonization of heterogeneous vertical navigation systems. Historically, global aviation has relied on a dual-reference architecture: manned aviation utilizes Barometric Pressure Altitude referenced to the International Standard Atmosphere (ISA) to ensure relative separation within the same air mass, whereas UAS predominantly operate on Geometric Altitude (Mean Sea Level, MSL) derived from Global Navigation Satellite Systems (GNSS) \cite{simonetti2024geodetic, zhao2023vertical}.

These two measurement domains—geopotential and geometric—are physically distinct. Due to non-standard atmospheric conditions (e.g., temperature deviations from ISA, pressure gradients) and geoid undulations, the discrepancy between pressure altitude and true geometric height can exceed hundreds of meters \cite{narayanan2022enhanced}. In obstacle-rich urban environments where UAS operate with minimal vertical safety margins, such discrepancies pose unacceptable collision risks and severely hinder the deployment of systematic U-space services.

% \subsection{Current Methodologies and Limitations}

Bridging the gap between barometric and geometric altitude has evolved from early inertial-barometric integrations to modern data-driven methodologies. State-of-the-art approaches leverage high-fidelity Numerical Weather Prediction (NWP) models, such as ECMWF's ERA5, as "virtual sensors" to reconstruct the atmospheric profile and correct the barometric formula \cite{schumann2017use}. For instance, algorithms like Baro-ARAIM and BiG-C have demonstrated the capability to reduce vertical errors to approximately 4--30 meters by dynamically estimating temperature lapse rates and pressure fields \cite{simonetti2024geodetic, narayanan2022enhanced}.

Despite these advancements, existing methodologies exhibit critical limitations in urban VLL airspace:
\begin{enumerate}
    \item \textbf{Spatiotemporal Latency}: High-fidelity NWP models often suffer from coarse grid resolutions (e.g., $0.25^\circ$) and publication latency, failing to capture high-frequency microclimate variations (e.g., urban heat islands, building wake turbulence) inherent to urban canyons.
    \item \textbf{Sensor-Specific Biases}: Low-cost micro-electro-mechanical systems (MEMS) barometers deployed in commercial UAS and edge IoT nodes are highly susceptible to environmental influences, thermal drift, and aerodynamic interference, complicating absolute height determination \cite{sun20253d, li2023precise}.
    \item \textbf{Poor Spatial Generalization}: Traditional machine learning approaches (e.g., Random Forests) trained to correct these localized errors typically act as "black boxes." They often violate physical atmospheric laws in data-sparse regions, leading to severe extrapolation errors when deployed to unseen urban topologies \cite{zhang2020urban}.
\end{enumerate}

% \subsection{Proposed System and Contributions}

To address these instrumentation and measurement challenges, this paper proposes a novel multi-sensor fusion framework comprising a low-cost edge-sensing terminal (the "Height Box") and a Physics-Informed Neural Field (PINF) algorithm. Instead of solely relying on delayed global weather models or unconstrained machine learning, our system deploys distributed edge instruments to capture synchronized 4D atmospheric profiles. We then formulate the altitude conversion as a Physics-Informed Zero-Shot Generalization framework, utilizing physics-derived pressure bias to standardize inputs across uncalibrated sensors. By embedding the hydrostatic equation into the neural architecture, predicting pressure corrections rather than direct heights, and utilizing multi-resolution hash encoding, the proposed system learns high-frequency urban spatial features while strictly adhering to physical atmospheric constraints.

The main contributions of this work are as follows:
\begin{enumerate}
    \item \textbf{An Edge-Sensing Measurement Framework}: We present the design and deployment of a multi-sensor acquisition system that fuses GNSS, barometric, and hygrometric data, providing a scalable hardware solution for urban atmospheric sampling. Crucially, the system natively resolves real-time Mean Sea Level (MSL) altitude by outputting orthometric heights coupled with geoid anomaly offsets, bridging the geometric and barometric domains directly at the edge without requiring computationally expensive offline geoid modeling.
    \item \textbf{Physics-Informed Zero-Shot Generalization}: We pioneer the use of Neural Fields predicting environmental pressure corrections $\delta P$ rather than arbitrary height residuals. By decoupling sensor-specific errors through a physics-derived pressure bias, our network generalizes universally across completely unseen hardware.
    \item \textbf{Altitude-Based Curriculum Learning}: To overcome the challenge of atmospheric gradient non-linearities and spatial heterogeneity, we introduce a novel three-stage altitude-based curriculum learning strategy. Progressing from simple surface patterns ($<100\text{m}$) to extreme outlier cases, it accelerates convergence and fundamentally increases stability.
    \item \textbf{Rigorous Metrological Validation}: Unlike prior studies relying on random data splits that inadvertently leak sensor bias across folds, we validate our system using a strict 8-fold Leave-One-Sensor-Out (LOSO) testing protocol. Achieving a true zero-shot Mean Absolute Error (MAE) of 3.55 meters, we demonstrate the robustness and operational capability required for heterogeneous low-altitude traffic.
\end{enumerate}

\section{Related Work}
\label{sec:related_work}

\subsection{Analytical and NWP-Based Altitude Harmonization}
Vertical navigation has historically operated on a dual-reference system. Manned aviation utilizes pressure altitude, which references a virtual datum (ISA MSL) to ensure relative separation between aircraft within the same air mass \cite{SIDs_2022}. In contrast, GNSS provides geometric altitude relative to the WGS-84 ellipsoid. The transformation between these systems involves correcting for two primary factors: the separation between the geoid and the ellipsoid (Geoid Undulation, $N$), and the deviation of the actual atmospheric temperature profile from the ISA model (Scale Factor Error) \cite{navi_637}. The ICARUS project highlighted that without real-time harmonization services, the mixing of barometric and geometric traffic in Very Low-Level (VLL) airspace could lead to catastrophic vertical separation infringements \cite{SIDs_2022}.

To bridge this gap, classic methods, such as the Blanchard algorithm~\cite{blanchard1971new}, integrate inertial measurements with air data to estimate errors induced by non-standard temperature and pressure gradients. However, these methods rely heavily on the accuracy of onboard temperature sensors and are prone to drift over long durations \cite{li2010errors}. More recently, researchers have leveraged Numerical Weather Prediction (NWP) data to act as a "virtual sensor". Simonetti and Garc\'{i}a Crespillo demonstrated that interpolating 4D weather data (e.g., ERA5) can reduce barometric altitude errors to meter-level~\cite{navi_637}. Similarly, Schumann \textit{et al.} utilized NWP analysis to verify aircraft height-keeping performance, proving that meteorological reanalysis can serve as a truth source for static pressure calibration \cite{Schumann_2017}. Despite their high accuracy, NWP-based methods are limited by the latency of data availability and the coarse spatial resolution of global models (typically $0.25^\circ$), which frequently fail to capture microclimate variations in urban canyons.

\subsection{Data-Driven Calibration for Barometric Sensors}
Conventional methods for barometric altitude estimation rely on the barometric formula with empirical corrections \cite{icao1993manual}. With the advent of deep learning, data-driven algorithms have been widely applied to model the non-linear relationship between pressure and geometric height. Recent work incorporates machine learning to learn sensor-specific biases, with ensemble approaches like Random Forests typically achieving 10 $\sim$ 30m accuracy in urban altitude estimation \cite{zhang2020urban}.

While standard neural networks and tree-based models can approximate these non-linear mappings, they act as ``black boxes'' and lack physical interpretability. Consequently, they often lead to poor generalization outside the training distribution (extrapolation), as they do not inherently respect the hydrostatic constraints governing atmospheric physics. Furthermore, these standard approaches traditionally lack explicit architectural mechanisms to decouple static unit hardware bias from the environmental model, inherently preventing zero-shot generalization across previously unseen, uncalibrated instruments in real-world deployments.

\subsection{Neural Fields and Spatial Representation}
To achieve continuous and robust spatial modeling, Coordinate-Based Neural Networks, or Neural Fields, have emerged as a powerful paradigm. Originally popularized by Neural Radiance Fields (NeRF) for novel view synthesis \cite{mildenhall2020nerf}, this architecture maps spatial coordinates to continuous physical quantities. To overcome the spectral bias of standard Multi-Layer Perceptrons (MLPs), M\"{u}ller \textit{et al.} introduced Instant-NGP, utilizing a multi-resolution hash encoding technique that dramatically accelerates the learning of high-frequency spatial details \cite{muller2022instant}.

While Neural Fields have revolutionized computer vision, their application to continuous geospatial regression and massive atmospheric parameter estimation remains largely underexplored. Furthermore, training Neural Fields on spatially heterogeneous multi-sensor IoT data presents severe convergence methodologies. Inspired by Bengio \textit{et al.} \cite{bengio2009curriculum}, Curriculum Learning—a strategy of training models on progressively harder examples—offers a potential solution. In this paper, we pioneer the adaptation of Multi-Resolution Hash Encoding, intertwined with SIREN features, for high-altitude metrological applications. By coupling an altitude-based curriculum learning scheme with explicit formulation of physical hardware pressure bias, we transform continuous Neural Fields into a highly accurate, zero-shot generalizable spatial measurement topology tool.


% \section{Preliminary}
\label{sec:preliminary}
\subsection{Geodetic and Orthometric Height}
First, we clarify the geometric relationship between different vertical datums. The GNSS module measures the Geometric Altitude ($h$), defined as the height above the WGS-84 reference ellipsoid1. However, aviation and barometric systems operate relative to the Geoid (approximated by Mean Sea Level, MSL). The Orthometric Height ($H$) is related to $h$ by the Geoid Undulation ($N(\phi, \lambda)$):
$$H \approx h - N(\phi, \lambda) \ , $$
where $(\phi, \lambda)$ represents latitude and longitude. The undulation $N$ can vary significantly (e.g., -105m to +85m globally). In our framework, $N$ is retrieved from the EGM96/EGM2008 gravity model to align the geometric and geopotential frames.

\subsection{The Hydrostatic Atmosphere Model}
The fundamental physical law governing the relationship between pressure $P$ and geopotential height $H$ is the hydrostatic equation:$$\frac{dP}{dH} = -\rho g$$By incorporating the ideal gas law $\rho = \frac{P}{R_\text{dry} T_\text{v}}$, where $R_\text{dry}$ is the specific gas constant for dry air and $T_\text{v}$ is the virtual temperature (correcting for humidity effects), we derive the Hypsometric Equation:
$$H(P) = H_\text{ref} + \frac{R_\text{dry} \bar{T}_\text{v}}{g_0} \ln\left(\frac{P_\text{ref}}{P}\right) \ ,$$
where $P_\text{ref}$ and $H_\text{ref}$ are values at a reference level (e.g., ground level).

\subsection{The Residual Problem}
While NWP models (like ERA5) provide baseline estimations for $T_\text{v}$ and $P_\text{ref}$, they fail to capture micro-scale disturbances in urban canyons (e.g., heat islands, building wake turbulence) due to coarse grid resolution (>20km). We define the Height Residual ($R_{\Delta}$) as the discrepancy between the GNSS-measured true height and the physics-derived height:$$h_{GNSS} = \mathcal{F}_{phy}(P_{obs}, T_{obs}, \text{NWP}) + R_{\Delta}(x, y, z, t)$$The goal of our methodology is to model this continuous residual field $R_{\Delta}$.

\section{Methodology}
\label{sec:methodology}

The objective of this research is to establish a rigorous, high-precision conversion framework between barometric pressure and geometric altitude in complex urban environments. Unlike pure black-box deep learning approaches, our methodology decomposes the altitude measurement into a deterministic physical baseline and a learnable spatial residual component. To achieve this, we co-design an edge-sensing hardware terminal and a Physics-Informed Neural Field (PINF) algorithm.

\begin{figure}[htbp]
    \centering
    % Name your image fig2_hardware.png and put it in the image folder
    \includegraphics[width=\columnwidth]{image/figure2.pdf}
    \caption{Hardware architecture and signal flow of the GeoBox data acquisition terminal. The system integrates GNSS, barometric, and environmental sensors. An onboard MCU performs strict temporal synchronization and simple sensor fusion before transmitting the data payload via a 4G LTE/MQTT interface.}
    \label{fig:hardware}
\end{figure}

\subsection{Edge Sensing Instrumentation}
\label{subsec:hardware}

The foundation of our proposed framework is a custom-designed, low-cost multi-sensor data acquisition terminal, designated as the ``GeoBox''. This edge device is engineered to capture synchronized navigation and atmospheric profiles across distributed urban topologies, providing the sparse but high-fidelity data required to train the continuous neural field. The hardware architecture integrates three core sensing modules:
\begin{enumerate}
    \item \textbf{GNSS Positioning Module}: An integrated GNSS receiver coupled with an external active antenna acquires high-precision geometric coordinates (latitude $\phi$, longitude $\lambda$) and the true Geometric Altitude ($h_\text{GNSS}$) relative to the WGS-84 reference ellipsoid. It supports multiple constellations (GPS, BeiDou, Galileo) to ensure lock stability in urban canyons.
    \item \textbf{Barometric Pressure Sensor}: A high-resolution MEMS-based static pressure sensor provides raw atmospheric pressure measurements ($P_\text{obs}$) with an operating range of 300 to 1100 hPa. The sensor is isolated within a specially designed enclosure featuring vented static ports to minimize dynamic pressure artifacts induced by wind gusts or platform motion.
    \item \textbf{Environmental Sensing Array}: A digital hygrometer and thermometer are integrated to measure ambient temperature ($T_\text{obs}$) and relative humidity ($RH$). These parameters are strictly required to calculate the specific humidity and virtual temperature, effectively compensating for air density variations in the subsequent hydrostatic calculations.
\end{enumerate}

\subsubsection{Data Synchronization and Transmission}
Metrological integrity depends heavily on the temporal synchronization of the heterogeneous data streams. An onboard Microcontroller Unit (MCU) performs time synchronization, ensuring that the GNSS geometric altitude is strictly time-stamped alongside the barometric and environmental readings. To guarantee data reliability during field deployments, the system executes a Built-In Test (BIT) sequence upon startup to verify GNSS lock status, network connectivity, and sensor health. Powered by a lithium-ion battery for extended autonomy, the MCU aggregates the data into a standardized payload—comprising a unique device ID (UID), Unix timestamp, $P_\text{obs}$, $T_\text{obs}$, $RH$, and $(\phi, \lambda)$ coordinates. This packet is transmitted in real-time to a centralized cloud server via a 4G/LTE cellular module utilizing the MQTT protocol. This continuous, synchronized data pipeline forms the empirical basis for the physical baseline estimation and neural residual learning phases.

\subsection{Physical Baseline Estimation and Problem Formulation}
\label{subsec:physical_baseline}

To isolate the complex, non-linear microclimate disturbances inherent to urban VLL airspace, we first establish a deterministic physical baseline. This baseline serves as the mean function for our estimation framework, derived from fundamental atmospheric physics and the edge-sensor measurements acquired by the GeoBox.

The foundational relationship between atmospheric pressure $P$ and geopotential height $H$ is governed by the hydrostatic equation:
\begin{equation}
\frac{dP}{dH} = -\rho g \ ,
\label{eq:hydrostatic}
\end{equation}
where $\rho$ is the air density and $g$ is the gravitational acceleration. Standard aviation models (e.g., ISA) assume a static dry-air atmosphere \cite{icao1993manual}. However, to achieve meter-level accuracy, variations in air density caused by moisture must be accounted for. We utilize the onboard environmental sensors to compute the in-situ Virtual Temperature $T_\text{v}$:
\begin{equation}
T_\text{v} = T_\text{obs} (1 + 0.61 q) \ ,
\label{eq:virtual_temp}
\end{equation}
where $T_\text{obs}$ is the measured absolute temperature and $q$ is the specific humidity derived from the measured relative humidity ($RH$) \cite{simonetti2024geodetic}.

By substituting the ideal gas law ($\rho = P / (R_\text{dry} T_\text{v})$) into Eq. (\ref{eq:hydrostatic}) and integrating, we obtain the Hypsometric Equation, which defines our physical baseline estimator $\hat{H}_\text{phy}$:
\begin{equation}
\hat{H}_\text{phy} = H_\text{ref} + \frac{R_\text{dry} T_\text{v}}{g_0} \ln\left(\frac{P_\text{ref}}{P_\text{obs}}\right) \ ,
\label{eq:hypsometric}
\end{equation}
where $R_\text{dry}$ is the specific gas constant for dry air, $g_0$ is the standard gravity, and $P_\text{ref}$ and $H_\text{ref}$ are the reference pressure and height, respectively. This step accounts for the primary altitude variance driven by macro-meteorological conditions.

Crucially, $\hat{H}_\text{phy}$ represents the orthometric height (relative to the global mean sea level, MSL). To harmonize this standard atmospheric altitude with standard UAS GNSS coordinate frames—which typically output Height Above Ellipsoid (HAE)—one must traditionally apply the Geoid Undulation anomaly $N(\phi, \lambda)$ retrieved from high-resolution gravity models (e.g., EGM96/EGM2008).
However, deploying global high-resolution geoid models heavily limits real-time edge processing capabilities. In our proposed GeoBox edge-sensing architecture, the integrated multi-constellation GNSS receiver is uniquely engineered to output true orthometric heights ($h_\text{GNSS} \equiv$ MSL) directly by utilizing onboard regional geoid anomaly grids. Therefore, we utilize the mathematically rigorous $\hat{H}_\text{phy}$ directly as the un-deformed physical baseline $h_\text{phy}$, decoupling the high-frequency measurement calibration from expensive offline gravitational geoid interpolations.

While $h_\text{phy}$ provides a strong generalized mean function, it fundamentally fails to capture high-frequency local atmospheric disturbances, such as urban heat islands, building wake turbulence, and importantly, the individual static bias inherent to uncalibrated hardware, which cannot be modeled by simple analytical equations. To enforce physical consistency across multiple uncalibrated nodes, we introduce a Physics-Derived Pressure Bias Formulation. Instead of directly predicting an explicit geometric residual ($\Delta h$), we predict the localized barometric pressure correction ($\delta P$).

The unified urban altitude estimation problem is thus formulated as calculating a pressure correction field and applying it to the physical hydrostatic reconstruction:
\begin{equation}
P_\text{corr} = P_\text{obs} + f_\theta(\mathbf{x}, \mathbf{p}, t) \ ,
\label{eq:pressure_correction}
\end{equation}
where $f_\theta$ represents the proposed PINF parameterized by weights $\theta$. The Neural Field predicts the continuous pressure residual $\delta P$ using spatiotemporal coordinates $\mathbf{x} = (\phi, \lambda, z)$, temporal mapping $t$, and physical features $\mathbf{p}$ (importantly including the sensor pressure bias). The final predicted height $h_\text{pred}$ is strictly routed through the un-deformed Hypsometric Equation using the corrected pressure $P_\text{corr}$:
\begin{equation}
h_\text{pred} = H_\text{ref} + \frac{R_\text{dry} T_\text{v}}{g_0} \ln\left(\frac{P_\text{ref}}{P_\text{corr}}\right) \ .
\label{eq:problem_formulation}
\end{equation}

\begin{figure}[htbp]
    \centering
    % Ensure your assembled image is named fig3_framework.png and placed in the image folder
    \includegraphics[width=1\linewidth]{image/fig3.pdf}
    \caption{The overall algorithmic framework. The Physical Baseline Branch derives a deterministic geometric height ($\hat{h}_{phy}$) based on the hypsometric equation and geoid undulation. Concurrently, the Neural Residual Branch leverages multi-resolution hash encoding and engineered terrain features to learn the complex spatial altitude residual ($\Delta h$). The final prediction ($h_{pred}$) is supervised by GNSS ground truth, with backpropagation (dashed line) restricted to the neural weights to preserve physical interpretability.}
    \label{fig:framework}
\end{figure}

\subsection{Spatiotemporal Encoding and Physics Bias Formulation}
\label{subsec:hash_and_terrain}

To effectively learn the highly non-linear residual field $\delta P$, the neural network must be capable of resolving high-frequency spatial variations induced by urban microclimates without heavily overfitting to single sensors. Standard coordinate-based Multi-Layer Perceptrons (MLPs) inherently suffer from ``spectral bias'', favoring low-frequency functions and struggling to fit sharp spatial gradients. To overcome this, we adapt the multi-resolution hash encoding technique~\cite{muller2022instant} alongside a novel temporal Fourier encoding map.

\subsubsection{Multi-Resolution Spatiotemporal Encoding}
Unlike traditional positional encoding that relies on fixed sinusoidal functions, hash encoding maps the spatial coordinates $(x, y, z)$ to trainable feature vectors across multiple resolution scales. For a given resolution level $l \in \{1, 2, \dots, L\}$ with a corresponding grid resolution $N_l$, we define a hash table $\mathcal{H}_l$ that maps spatial indices to continuous feature vectors:
\begin{equation}
\mathbf{e}_l(x, y, z) = \mathcal{H}_l\left(\text{hash}\left(\lfloor x \cdot N_l \rfloor, \lfloor y \cdot N_l \rfloor, \lfloor z \cdot N_l \rfloor\right)\right)
\label{eq:hash_lookup}
\end{equation}
where $\mathbf{e}_l \in \mathbb{R}^F$ is the feature vector with dimension $F=4$, and we utilize $L=16$ independent levels.

The grid resolutions follow a geometric progression to ensure a balanced receptive field:
\begin{equation}
N_l = N_{\text{min}} \cdot b^l, \quad b = \left(\frac{N_{\text{max}}}{N_{\text{min}}}\right)^{\frac{1}{L-1}}
\label{eq:resolution}
\end{equation}
with $N_{\text{min}} = 16$ and $N_{\text{max}} = 512$. The final spatial encoding $\phi(x, y, z)$ is the concatenation of features interpolated from all $L$ levels:
\begin{equation}
\phi(x, y, z) = [\mathbf{e}_1, \mathbf{e}_2, \dots, \mathbf{e}_L] \in \mathbb{R}^{L \times F}
\label{eq:concat_encoding}
\end{equation}
Furthermore, to model periodic atmospheric variations, the temporal observation variable $t$ is projected into a high-dimensional space via Fourier transformations across 6 distinct frequencies.

\subsubsection{Decoupling Hardware via Physics-Derived Pressure Bias}
The single greatest challenge for deep-learning-based metrology is generalizing to completely unseen uncalibrated sensor nodes. Direct MLPs will inherently memorize the systematic hardware bias of training units. To enforce zero-shot generalization and abstract the network away from strict node identities, we propose a Physics-Derived Pressure Bias formulation point:
\begin{equation}
P_\text{bias} = P_\text{obs} - P_\text{expected}
\label{eq:bias_formulation}
\end{equation}
where $P_\text{expected}$ is the ideal barometric reading analytically derived by inversing the uncalibrated GNSS altitude ($h_\text{GNSS}$) against the deterministic Hypsometric physical baseline.

By calculating $P_\text{bias}$ and feeding it as a learned embedding rather than using discrete sensor-ID tokens or vulnerable local terrain densities, the network fundamentally acts on the relative discrepancy space. To predict the pressure correction $\delta P$, this explicit physics characteristic is concatenated with the spatial hash encoding, temporal Fourier encoding, temperature $T$, and relative humidity $RH$ to form the complete continuous feature vector:
\begin{equation}
\mathbf{f}_{\text{input}} = [\phi(x, y, z), \phi_{\text{Fourier}}(t), P_\text{bias}, T, RH]
\label{eq:full_input}
\end{equation}
This fused vector $\mathbf{f}_{\text{input}}$ is subsequently processed by the Neural Residual Field's MLP to output the highly precise pressure correction $\delta P$.

\subsection{Altitude-Based Curriculum Learning Strategy}
\label{subsec:curriculum_learning}

Training continuous Neural Fields on massive and spatially heterogeneous IoT data presents severe convergence challenges. In an urban VLL environment, measurements acquired near ground-level arrays exhibit strong spatial correlation and are less perturbed by altitude-dependent atmospheric degradation, making them mathematically ``easier'' to fit. Conversely, measurements extending heavily vertically into the atmosphere are sparse, boundary-defining constraints prone to non-linear physical drift, representing ``harder'' samples. To algorithmically prevent the network from overfitting to densely sampled ground-level data while failing at higher altitudes, we propose a strict Altitude-Based Curriculum Learning structure.

\subsubsection{Three-Stage Curriculum Scheme}
Inspired by cognitive learning principles \cite{bengio2009curriculum}, the training process is explicitly structured to expose the SIREN-based geometry network to progressively more challenging geospatial distributions. We dynamically partition the full training dataset $\mathcal{D}$ into three distinct curriculum stages bound by absolute geometric altitude ($h$):
\begin{itemize}
    \item \textbf{Stage 1 (Easy - Foundation Formulation)}: The network is exclusively trained on low-altitude samples yielding stable thermodynamic behaviors. $\mathcal{D}_1 = \{ \mathbf{x} \in \mathcal{D} \mid h < 100\text{m} \}$. This safely establishes fundamental spatial hash structures.
    \item \textbf{Stage 2 (Medium - Moderate Expansion)}: The distribution bounds are relaxed to introduce mild-altitude geometries and early extrapolation behaviors. $\mathcal{D}_2 = \{ \mathbf{x} \in \mathcal{D} \mid h < 200\text{m} \}$.
    \item \textbf{Stage 3 (Hard - Extreme Outliers)}: The network is finally optimized against the entire dataset, introducing severe, high-altitude outliers. $\mathcal{D}_3 = \mathcal{D}$.
\end{itemize}

During training, the model systematically transitions from $\mathcal{D}_1 \rightarrow \mathcal{D}_2 \rightarrow \mathcal{D}_3$. Stage 1 and 2 operate for 30 epochs each, constructing stable global gradients, before Stage 3 fine-tunes on the complete unconstrained atmospheric volume for 80 epochs. If directly exposed to extreme vertical disparities without foundation (No Curriculum Strategy), gradient explosion fundamentally limits valid generalization to 9.27m MAE. By employing the 3-stage altitude scheme, the optimized framework confidently stabilizes across all bounds.

\begin{figure}
    \centering
    \includegraphics[width=1\linewidth]{image/fig5_curriculum.png}
    \caption{Curriculum learning strategy effectively preventing gradient collapse and stabilizing spatial hash geometries by sequencing altitude complexity.}
    \label{fig:curriculum}
\end{figure}

\subsubsection{SIREN Architecture and Optimization}
The fused feature vector $\mathbf{f}_{\text{input}}$ is linearly projected through a 3-layer explicit Multi-Layer Perceptron (hidden dimension=256) to output the localized $\delta P$. Crucially, instead of standard ReLU/SiLU components, we utilize Sinusoidal Representation Networks (SIREN) for the activation layers. SIREN architectures inherently preserve analytical derivative stability for periodic physical signals, preventing vanishing gradient stagnation over harsh geographical boundaries.

The forward pass is evaluated continuously per batch using an AdamW optimizer iteratively minimizing the exact residual between reconstructed orthometric height and strict observed GNSS MSL:
\begin{equation}
\mathcal{L}(\theta) = \frac{1}{B} \sum_{i=1}^B \left| h^\text{pred}_{i} - h^{true}_i \right|
\label{eq:loss_function}
\end{equation}
Furthermore, a Cosine Annealing learning rate scheduler with warm restarts is initialized ($10^{-3}$), effectively allowing periodic deep landscape exploration and preventing premature optimization stagnation.

% \subsubsection{Training Efficiency}
% \subsubsection{Future Directions}

% \begin{enumerate}
%     \item \textbf{Meta-learning}: Rapid adaptation to new sensors with few samples
%     \item \textbf{Physics-informed neural operators}: Incorporate PDE constraints for better extrapolation
%     \item \textbf{3D atmospheric fields}: Estimate full 3D pressure fields rather than single altitudes
%     \item \textbf{Transfer learning}: Apply trained models to new cities with minimal calibration
% \end{enumerate}

%%%%%%%%%%%%%%%%%%%%%%%%%%%%%%%%%%%%%%%%%%%%%%%%%%%%%%%%%%%%%%%%%%%%%%%%%%%%%%%%%%%%%%%%%%%%%%%%%%%%%%%%%%%%%%%%%%%%%%%%%%%%%%%%%%%%%%%%%%%%%%%%%%%%%%%%%%%%



\section{Experiments and Results}
\label{sec:experiments}

To rigorously validate the proposed multi-sensor fusion framework, an extensive urban field campaign was conducted. This section details the hardware deployment of the designed Geobox and a comprehensive metrological evaluation of the altitude measurement performance against state-of-the-art baselines.

\subsection{Experimental Setup and Hardware Deployment}
\label{subsec:exp_setup}

The field experiments were conducted in a densely built urban sector in Shenzhen, China, covering approximately $1 \text{ km}^2$. This specific area is characterized by complex urban canyons, providing a highly challenging, real-world metrological testbed for barometric altimetry. We deployed a network of our custom-designed Geobox across various urban topologies, including ground-level infrastructure and mid-rise building rooftops. We record the raw IoT data streams for about two weeks, which generated a high-fidelity dataset comprised $N = 115,417$ synchronized samples. Table~\ref{tab:dataset} summarizes dataset statistics.

To supply the necessary meteorological priors, ERA5 atmospheric reanalysis data from the ECMWF (including 2-meter temperature $t_{2m}$ and surface pressure $sp$ at $0.25^\circ$ resolution) was interpolated and integrated into the dataset. Finally, a global physical calibration established the baseline barometric parameters: scale height $H_s = 9690.68 \text{ m}$ and reference pressure $P_0 = 101698.84 \text{ Pa}$.

The training setup is as follows: nvidia L20 GPU is used. Batch size is xxx. 

\begin{table}[t]
\caption{Dataset Statistics}
\label{tab:dataset}
\centering
\begin{tabular}{lc}
\toprule
\textbf{Statistic} & \textbf{Value} \\
\midrule
Total samples & 115,417 \\
Number of sensors & 7 \\
Spatial coverage & $\sim$1 km$^2$ \\
Sampling frequency & 1 minute \\
\midrule
\multicolumn{2}{c}{\textit{Sensor Heights (m)}} \\
\midrule
Sensor 42499896 & 58.0 \\
Sensor 78250224 & 139.0 \\
Sensor 42508217 & 100.1 \\
Sensor 27528610 & 95.3 \\
Sensor 27536362 & 107.9 \\
Sensor 78251938 & 95.8 \\
Sensor 27373510 & 259.0 \\
\bottomrule
\end{tabular}
\end{table}


\subsection{Metrological Evaluation Protocol}
\label{subsec:eval_protocol}

Evaluating spatial altitude models using random data splits inevitably leads to data leakage and overestimates performance due to temporal and spatial autocorrelation. To rigorously validate the system's metrological reliability and its capability to generalize to unseen urban coordinates, we employ a strict \textbf{Leave-One-Sensor-Out (LOSO) cross-validation} protocol. In each cross-validation fold, the neural field is trained exclusively on the spatial domain of $K-1$ sensors and evaluated entirely on the unseen geographic location of the held-out sensor. To contextualize the performance, our framework is benchmarked against four distinct baselines:
\begin{enumerate}
    \item \textbf{Physics Baseline}: The deterministic hypsometric equation without neural residual compensation.
    \item \textbf{Random Forest}: A classical ensemble decision tree method representing traditional data-driven sensor calibration \cite{zhang2020urban}.
    \item \textbf{Basic Neural Field}: A standard Multi-Layer Perceptron (MLP) utilizing simple positional encoding.
    \item \textbf{NF + ERA5}: A Neural Field directly incorporating ERA5 meteorological features without multi-resolution hashing.
    % \item \textbf{SIREN + Ensemble}: The prior state-of-the-art spatial modeling approach utilizing periodic sinusoidal activation functions and model ensembling \cite{sitzmann2020implicit}.
\end{enumerate}
The primary evaluation metrics are the Mean Absolute Error (MAE) and Root Mean Squared Error (RMSE) against the GNSS-derived ground truth, which are defined as:
\begin{itemize}
    \item \textbf{Mean Absolute Error (MAE)}: $\frac{1}{N} \sum_{i=1}^N |h_i^{\text{pred}} - h_i^{\text{true}}|$.
    \item \textbf{Root Mean Squared Error (RMSE)}: $\sqrt{\frac{1}{N} \sum_{i=1}^N (h_i^{\text{pred}} - h_i^{\text{true}})^2}$.
    \item \textbf{Improvement over baseline}: $\frac{\text{MAE}_{\text{baseline}} - \text{MAE}_{\text{method}}}{\text{MAE}_{\text{baseline}}} \times 100\%$.
\end{itemize}




\subsection{Altitude Measurement Performance}
\label{subsec:main_results}

The nominal altitude measurement performance of the proposed PINF framework, alongside the five baseline methodologies, is quantitatively summarized in Table \ref{tab:main_results}. The results demonstrate that our framework significantly outperforms both traditional analytical methods and contemporary data-driven architectures in complex urban environments. As expected, the standard \textbf{Physics Baseline} exhibits the highest error (MAE = 35.03 m). While the hypsometric equation accounts for macro-meteorological variables, it is fundamentally incapable of resolving the high-frequency microclimate disturbances inherent to urban canyons. The classical \textbf{Random Forest} calibrator reduces the error to 22.00 m. However, its performance under the strict LOSO protocol reveals a critical limitation of traditional tree-based models: severe spatial overfitting. They fail to mathematically extrapolate to unseen sensor coordinates because they lack underlying physical constraints. Among the neural architectures, our proposed \textbf{PINF framework} drastically suppresses the residual error, achieving a remarkable \textbf{MAE of 3.79 m} and an RMSE of 4.85 m. This represents an \textbf{89.2\% improvement} over the uncalibrated physics baseline. By bringing the absolute vertical error consistently below the 4-meter threshold, the PINF framework satisfies the stringent vertical separation requirements for high-density U-space operations, proving the viability of using low-cost edge IoT devices for precision metrology.

\begin{table}[htbp]
\caption{Metrological Performance Comparison on Urban Altitude Estimation (LOSO Validation)}
\label{tab:main_results}
\centering
\begin{tabular}{lccc}
\toprule
\textbf{Method} & \textbf{MAE (m)} & \textbf{RMSE (m)} & \textbf{Improv. (vs. Phys.)} \\
\midrule
Physics Baseline & 35.03 & 42.15 & - \\
Random Forest \cite{zhang2020urban} & 22.00 & 28.50 & 37.2\% \\
Basic Neural Field & 16.66 & 22.10 & 52.4\% \\
NF + ERA5 & 14.13 & 18.45 & 59.7\% \\
\midrule
\textbf{Proposed} & \textbf{3.79} & \textbf{4.85} & \textbf{89.2\%} \\
\bottomrule
\end{tabular}
\end{table}


\subsection{Ablation Study and Component Analysis}
\label{subsec:ablation}

To isolate and quantify the metrological contribution of each engineered component within our framework, a comprehensive ablation study was conducted. We sequentially integrated advanced modules into a basic Neural Field, evaluating the corresponding reduction in measurement error. The ablation trajectory is detailed in Table \ref{tab:ablation}.

\begin{table}[htbp]
\caption{Ablation Study: Component-Wise Error Mitigation}
\label{tab:ablation}
\centering
\begin{tabular}{lcc}
\toprule
\textbf{Configuration} & \textbf{MAE (m)} & \textbf{Relative Error Gain} \\
\midrule
Basic Neural Field & 14.13 & - \\
\quad + ERA5 Integration & 11.19 & - 2.94 m ($\downarrow 20.8\%$) \\
\quad + SIREN Activation & 9.85 & - 1.34 m ($\downarrow 12.0\%$) \\
\midrule
\textbf{Our Core Upgrades:} & & \\
Multi-Res Hash Encoding & 6.42 & - 2.24 m ($\downarrow 25.9\%$) \\
\quad + Curriculum Learning & 4.85 & - 1.57 m ($\downarrow 24.5\%$) \\
\quad + Terrain Features & \textbf{3.79} & \textbf{- 1.06 m ($\downarrow 21.9\%$)} \\
\bottomrule
\end{tabular}
\end{table}


The analysis yields several critical insights into spatial altitude modeling:
\begin{itemize}
    \item \textbf{The Superiority of Hash Encoding:} Replacing global SIREN activations with \textbf{Multi-Resolution Hash Encoding} provided the single largest performance leap ($\downarrow 2.24$ m). Global MLPs inherently suffer from spectral bias, smoothing over the sharp pressure gradients found between adjacent urban structures. Hash encoding localizes feature learning, allowing the network to explicitly memorize and interpolate high-frequency microclimate pockets without degrading the global atmospheric trend.
    \item \textbf{Impact of Curriculum Learning:} The introduction of the \textbf{Terrain-Aware Curriculum Strategy} further reduced the MAE by 1.57 m. By explicitly forcing the network to converge on dense, well-behaved low-altitude data before exposing it to sparse boundary conditions, the curriculum mechanism effectively prevents the optimizer from settling into suboptimal local minima caused by spatial heterogeneity.
    \item \textbf{Physical Topology Constraints:} Finally, the integration of \textbf{Terrain Features} (roughness, height rank, and density) yielded a final 1.06 m reduction. This confirms our hypothesis that urban barometric anomalies are not purely stochastic; they are strongly correlated with the physical morphology of the surrounding infrastructure. Injecting these explicit geometric priors allows the PINF to better compensate for localized aerodynamic phenomena, such as building wake turbulence.
\end{itemize}



% \input{sections/conclusion}









% \begin{algorithm}[H]
% \caption{Weighted Tanimoto ELM.}\label{alg:alg1}
% \begin{algorithmic}
% \STATE 
% \STATE {\textsc{TRAIN}}$(\mathbf{X} \mathbf{T})$
% \STATE \hspace{0.5cm}$ \textbf{select randomly } W \subset \mathbf{X}  $
% \STATE \hspace{0.5cm}$ N_\mathbf{t} \gets | \{ i : \mathbf{t}_i = \mathbf{t} \} | $ \textbf{ for } $ \mathbf{t}= -1,+1 $
% \STATE \hspace{0.5cm}$ B_i \gets \sqrt{ \textsc{max}(N_{-1},N_{+1}) / N_{\mathbf{t}_i} } $ \textbf{ for } $ i = 1,...,N $
% \STATE \hspace{0.5cm}$ \hat{\mathbf{H}} \gets  B \cdot (\mathbf{X}^T\textbf{W})/( \mathbb{1}\mathbf{X} + \mathbb{1}\textbf{W} - \mathbf{X}^T\textbf{W} ) $
% \STATE \hspace{0.5cm}$ \beta \gets \left ( I/C + \hat{\mathbf{H}}^T\hat{\mathbf{H}} \right )^{-1}(\hat{\mathbf{H}}^T B\cdot \mathbf{T})  $
% \STATE \hspace{0.5cm}\textbf{return}  $\textbf{W},  \beta $
% \STATE 
% \STATE {\textsc{PREDICT}}$(\mathbf{X} )$
% \STATE \hspace{0.5cm}$ \mathbf{H} \gets  (\mathbf{X}^T\textbf{W} )/( \mathbb{1}\mathbf{X}  + \mathbb{1}\textbf{W}- \mathbf{X}^T\textbf{W}  ) $
% \STATE \hspace{0.5cm}\textbf{return}  $\textsc{sign}( \mathbf{H} \beta )$
% \end{algorithmic}
% \label{alg1}
% \end{algorithm}


\section{Conclusion and Future Work}
\label{sec:conclusion}

Accurate and unified vertical positioning in VLL airspace is a foundational prerequisite for the safety of emerging UAM operations. To bridge the critical measurement gap between barometric pressure altitude and GNSS geometric altitude, this paper proposed a novel edge-sensing and spatial modeling framework. 
On the instrumentation front, we designed and deployed the Geobox, a low-cost, multi-sensor IoT terminal capable of capturing synchronized, high-frequency atmospheric and geometric profiles in dense urban canyons. On the algorithmic front, we introduced a PINF architecture. By explicitly embedding the physical hypsometric equation as a deterministic baseline and utilizing multi-resolution hash encoding alongside engineered terrain features, the PINF effectively models the highly non-linear microclimate residuals that traditional analytical methods fail to capture. Furthermore, a Terrain-Aware Curriculum Learning strategy was implemented to ensure robust convergence across spatially heterogeneous data distributions. Evaluated through a rigorous LOSO metrological protocol on a massive real-world dataset collected in Shenzhen, China, the proposed framework achieved a MAE of 3.79 meters. This represents a substantial 89.2\% improvement over the uncalibrated physics baseline. The results conclusively demonstrate that coupling purpose-built edge IoT instrumentation with physically constrained neural architectures can reliably provide a meter-level unified height reference for heterogeneous U-space traffic.

% \subsection{ Future Work}
Future research will focus on investigating these high-altitude boundary effects. We aim to integrate computational fluid dynamics priors or real-time wind vector measurements into the neural field to explicitly compensate for aerodynamic pressure distortions. Additionally, we plan to explore meta-learning and few-shot adaptation techniques. This will enable the PINF framework to rapidly calibrate and generalize to entirely new urban topologies with minimal initial sensor deployment, further enhancing the scalability of the Height Box network for global UAM deployment.


% \section*{Acknowledgments}



% {
% \appendix[Proof of the Zonklar Equations]
% Use $\backslash${\tt{appendix}} if you have a single appendix:
% Do not use $\backslash${\tt{section}} anymore after $\backslash${\tt{appendix}}, only $\backslash${\tt{section*}}.
% If you have multiple appendixes use $\backslash${\tt{appendices}} then use $\backslash${\tt{section}} to start each appendix.
% You must declare a $\backslash${\tt{section}} before using any $\backslash${\tt{subsection}} or using $\backslash${\tt{label}} ($\backslash${\tt{appendices}} by itself
%  starts a section numbered zero.)
%  }

\bibliographystyle{IEEEtran}
\bibliography{ref}

\newpage
 
% \bf{If you include a photo:}\vspace{-33pt}
% \begin{IEEEbiography}[{\includegraphics[width=1in,height=1.25in,clip,keepaspectratio]{fig1}}]{Michael Shell}
% Use $\backslash${\tt{begin\{IEEEbiography\}}} and then for the 1st argument use $\backslash${\tt{includegraphics}} to declare and link the author photo.
% Use the author name as the 3rd argument followed by the biography text.
% \end{IEEEbiography}

% \vspace{11pt}

% \bf{If you will not include a photo:}\vspace{-33pt}
% \begin{IEEEbiographynophoto}{John Doe}
% Use $\backslash${\tt{begin\{IEEEbiographynophoto\}}} and the author name as the argument followed by the biography text.
% \end{IEEEbiographynophoto}

\vfill

\end{document}


